\section{Dolbeault moduli spaces}\label{section: dolbeault moduli spaces}

Recall tha the Dolbeault moduli space is the moduli space of semistable Higgs bundles, i.e.
\begin{gather}
\Higgs_n^d(X) \coloneqq \left< (\mc{E}, \varphi) \mymid \mc{E} \in \Bun_n^d(X), \; \varphi\colon \mc{E} \to \mc{E} \otimes \Omega_X \right>,\\
\mc{M}_{n, Dol}^d(X) \coloneqq \Higgs_n^{d, ss}(X).
\end{gather}
We can define similarly $\Higgs_{\underline{m}}^d(\mc{X})$ and $\mc{M}_{\underline{m}, Dol}^d(\mc{X})$.

We now want to extend \zcref{thm: eqv Bun} to the Dolbeault moduli spaces and prove the following.
\begin{thm} \label{thm: eqv Dol}
There is an equivalence of categories
\begin{equation} \label{equivalence dol}
\mc{M}_{n, Dol}^d(X) \cong \mc{M}_{n_{-d}, Dol}^0(\mc{X}).
\end{equation}
\end{thm}

As before, we will first extend \zcref{lemma Bun} to the moduli spaces of Higgs bundles and then \zcref{thm: eqv Dol} will follow. For convenience, we split the proof of \zcref{thm: eqv Dol} in intermediate lemmas.

\begin{lemma}
Under the isomorphism
\begin{equation}
\Omega_{\mc{X}} \cong \pi^* \Omega_X \otimes \O(\tfrac{1}{n}p)^{\otimes n-1}
\end{equation}
shown in \zcref{omega mc(X)}, the canonical morphism
\begin{equation} \label{eq: canonical morphism cotangent}
\pi^* \Omega_X \to \Omega_{\mc{X}} \cong \pi^*\Omega_X \otimes \O(\tfrac{n-1}{n}p)
\end{equation}
has the form $n \id_{\pi^* \Omega_X} \otimes s^{n-1}$, where $s$ is the tautological section of $\O(\tfrac{1}{n}p)$.
\end{lemma}

\begin{proof}
The coarse moduli map $\pi \colon \mc{X} \to X $ locally around $p$ has the following form
\[\begin{tikzcd}
	{\Ab^1 = \Spec k[x]} \\
	{\mc{X} = [\Ab^1 / \mu_n]} \\
	{X = \Ab^1 / \mu_n} & {\Spec k[x^n] = \Spec k[t].}
	\arrow["q'", from=1-1, to=2-1]
	\arrow["{x \mapsto t = x^n}"{pos=0.6}, "q = \pi \circ q'"' sloped, curve={height=-24pt}, from=1-1, to=3-2]
	\arrow["\pi", from=2-1, to=3-1]
	\arrow[between={0.4}{0.6}, equals, from=3-1, to=3-2]
\end{tikzcd}\]
Notice that we have the relation
\[
\dd t = \dd x^n = n \, x^{n-1} \, \dd x,
\]
therefore the canonical morphism $q^* \Omega_X \to \Omega_{\Ab^1}$ corresponds to
\begin{align}
	k[x] \otimes_{k[t]} k[t]\dd t \; &\to \; k[x]\dd x \\
	f(x) \otimes \dd t \quad &\mapsto \; f(x) \, n \, x^{n-1} \, \dd x.
\end{align}
Since $q'$ is étale, it induces isomorphisms on the cotangent sheaves. Moreover, the section $x$ corresponds via $q'$ to the tautological section $s$ of $\O(\tfrac{1}{n}p)$. In conclusion, we get that the canonical morphism $\pi^* \Omega_X \to \Omega_{\mc{X}} \cong \pi^*\Omega_X \otimes \O(\tfrac{n-1}{n}p)$ has the indicated form.
\end{proof}

Now we can use this morphism to extend the pullback functor $\pi^*$ to Higgs bundles. In particular, we need to lift the Higgs field.
Define
\begin{align}
	(\pi^*, \pi^{\#}) \colon \Higgs_n^d(X) &\to \Higgs_{n_0}^d(\mc{X}) \\
	(E, \varphi) &\mapsto (\pi^*E, \pi^{\#}\varphi),
\end{align}
where $\pi^{\#}\varphi$ is defined as the composition
\begin{multline}
\pi^{\#} \varphi \colon \pi^*E \xrightarrow{\pi^* \varphi} \pi^*(E \otimes \Omega_X) = \pi^*E \otimes \pi^* \Omega_X \\
\xrightarrow{\id \otimes n \hspace{0.3mm} \id \otimes s^{n-1}} \pi^*E \otimes \pi^*\Omega_X \otimes \O(\tfrac{n-1}{n}p) = \pi^*E \otimes \Omega_{\mc{X}}.
\end{multline}

Similarly, in the other direction define
\begin{align}
	(\pi_*, \pi_{\#}) \colon \Higgs_{n_0}^d(\mc{X})  &\to \Higgs_n^d(X)\\
	(\mc{E}, \eta) &\mapsto (\pi_*\mc{E}, \pi_{\#}\eta)
\end{align}
as follows.
From the equivalence \eqref{equivalence bun pt1}, it follows that $\mc{E} = \pi^*E$, where $E \cong \pi_* \mc{E}$. By definition, $\eta$ is a morphism in $\Hom_{\mc{X}}(\mc{E}, \mc{E}\otimes\Omega_{\mc{X}})$. Now, using \zcref{omega mc(X)}, we find a chain of natural isomorphisms
\begin{align}
\Hom_{\mc{X}}(\mc{E}, \, \mc{E}\otimes\Omega_{\mc{X}}) & \cong \Hom_{\mc{X}}\left( \pi^*E, \, \pi^*E \otimes \pi^*\Omega_X \otimes \O(\tfrac{n-1}{n}p) \right)\\
& \cong \Hom_X\left( E, \, \pi_*(\pi^*(E \otimes \Omega_{\mc{X}}) \otimes \O(\tfrac{n-1}{n}p)) \right) \label{eq: higgs field, adjunction} \\
& \cong \Hom_X\left( E, \, E \otimes \Omega_X \otimes \pi_*\O(\tfrac{n-1}{n}p) \right) \label{eq: higgs field, projection formula} \\
& \cong \Hom_X\left( E, \, E \otimes \Omega_X \right), \label{eq: higgs field, pushforward}
\end{align}
% \begin{gather}
% \Hom_{\mc{X}}(\mc{E}, \, \mc{E}\otimes\Omega_{\mc{X}}) \cong \Hom_{\mc{X}}\left( \pi^*E, \, \pi^*E \otimes \pi^*\Omega_X \otimes \O(\tfrac{n-1}{n}p) \right) \cong\\
% \Hom_X\left( E, \, \pi_*(\pi^*(E \otimes \Omega_{\mc{X}}) \otimes \O(\tfrac{n-1}{n}p)) \right) \cong\\
% \Hom_X\left( E, \, E \otimes \Omega_X \otimes \pi_*\O(\tfrac{n-1}{n}p) \right) \cong \\
% \Hom_X\left( E, \, E \otimes \Omega_X \right),
% \end{gather}
where \eqref{eq: higgs field, adjunction} follows by the adjunction $\pi^* \dashv \pi_*$, \eqref{eq: higgs field, projection formula} by the projection formula and \eqref{eq: higgs field, pushforward} by the isomorphism $\pi_* \O(\tfrac{n-1}{n}p) \cong \O_X$. \\
Composing the isomorphisms, we get
\begin{align}
	\Hom_X(E, E \otimes \Omega_X) &\xrightarrow{\cong} \Hom_{\mc{X}}(\mc{E}, \mc{E}\otimes \Omega_{\mc{X}}) \\
	\psi \; &\mapsto \; n \, \pi^* \psi \otimes s^{n-1}.
\end{align}
This shows that there exists a unique $\psi \in \Hom_X(E, E \otimes \Omega_X)$ such that $\eta = n \, \pi^* \psi \otimes s^{n-1}$, where in particular we have $\psi = \frac{1}{n} \, \pi_* \eta$. Therefore, we can define
\[
\pi_{\#}\eta \coloneqq \frac{1}{n} \, \pi_* \eta.
\]
This gives a well defined functor $(\pi_*, \pi_{\#})$.

We can now phrase the analogue of \zcref{lemma Bun}.

\begin{lemma}\label{lemma Dol}
The functors $(\pi^*, \pi^{\#})$ and $(\pi_*, \pi_{\#})$ define inverse equivalences of categories
\begin{equation}
\begin{tikzcd}
	{\Higgs_n^d(X)} & {\Higgs^d_{n_0}(\mc{X}).}
	\arrow[""{name=0, anchor=center, inner sep=0}, "{{(\pi^*, \pi^{\#})}}", curve={height=-18pt}, from=1-1, to=1-2]
	\arrow[""{name=1, anchor=center, inner sep=0}, "{{(\pi_*, \pi_{\#})}}", curve={height=-18pt}, from=1-2, to=1-1]
	\arrow["\cong"{description}, shift right=1.5, draw=none, from=0, to=1]
\end{tikzcd}
\end{equation}
\end{lemma}

\begin{proof}
At the level of vector bundles this is \zcref{lemma Bun}. We need to show that the same holds for the Higgs fields. Indeed, computing the two compositions we get
\begin{gather}
\varphi \xmapsto{\pi^{\#}} n \, \pi^* \varphi \otimes s^{n-1} \xmapsto{\pi_{\#}} \frac{1}{n} \, \pi_*(n \, \pi^* \varphi \otimes s^{n-1}) = \varphi\\
\eta \xmapsto{\pi_{\#}} \frac{1}{n} \, \pi_* \eta \xmapsto{\pi^{\#}} n \, \pi^*(\frac{1}{n} \, \pi_* \eta) \otimes s^{n-1} = n \, \pi^* \psi \otimes s^{n-1} = \eta.
\end{gather}
\end{proof}

As observed earlier, tensoring with a line bundle induces an equivalence, so we get
\[\begin{tikzcd}
	{\Higgs_n^d(X)} &&&& {\Higgs^0_{n_{-d}}(\mc{X}).}
	\arrow["{(\pi^*, \pi^{\#})(-) \otimes \O(-\tfrac{d}{n}p)}", shift left=2, from=1-1, to=1-5]
	\arrow["{(\pi_*, \pi^{\#})(\O(-\tfrac{d}{n}p) \otimes -)}", shift left=2, from=1-5, to=1-1]
\end{tikzcd}\]

Finally, we noticed that (semi)stability is preserved, therefore we get the desired equivalence
\begin{equation}\label{eq: equivalence of dolbeault moduli spaces}
\begin{tikzcd}
	{\mc{M}_{n, Dol}^d(X)} &&&& {\mc{M}^0_{n_{-d}, Dol}(\mc{X}).}
	\arrow["{(\pi^*, \pi^{\#})(-) \otimes \O(-\tfrac{d}{n}p)}", shift left=2, from=1-1, to=1-5]
	\arrow["{(\pi_*, \pi^{\#})(\O(-\tfrac{d}{n}p) \otimes -)}", shift left=2, from=1-5, to=1-1]
\end{tikzcd}
\end{equation}

This concludes the proof of \zcref{thm: eqv Dol}.